\subsection{Related work}\label{sec:related}

\paragraph{Single-call structured extraction.} Agrawal et al.\ showed that a general LLM, guided by a one-shot
example in the target format and followed by short deterministic \emph{resolvers}, could extract clinical
information such as medications and their status from English clinical snippets~\cite{agrawal2022fewshot}.
The same work reported two difficulties that bear directly on long forms: it remained ``difficult to guide a
LLM to match an exact schema'' when annotation guidelines span several pages, and prompts specialised for one
status tended to produce an answer even when none was present, a bias toward non-empty outputs. Wiest et al.\
ran Llama-2 models locally on 500 English history sections from MIMIC-IV to extract five binary features,
with an output format in which all features appear in one JSON object~\cite{privacy2024structured}. Unconstrained JSON was ``inconsistent and
defective'', so the authors enforced it with grammar-based sampling; the grammar also forced the model to
write an excerpt before each answer, but excerpts were not checked against the source and a hallucinated
excerpt was reported for the smallest model.

\paragraph{Specification and multi-step pipelines.} Hein et al.\ extracted diagnoses, specimen-level labels and
immunohistochemistry results from kidney-tumour pathology reports with a pipeline of segmentation prompts,
a standardisation prompt and a Python step that produces a long table, refined over six iterations with a
pathologist~\cite{iterativerefine2025}.\footnote{A preprint of the same study was posted under the title
\emph{Prompts to Table}; it reports the same data and results and is not counted as a separate study.} Their
error ontology separated model errors from annotation errors and from insufficient instructions (``schema
issues''); major errors were more frequent in reports with many tests, and no inter-annotator agreement was
measured because the reference was developed iteratively. The study illustrates that the specification of a
long target, not only the prompt, is an object of development and must be reported.

\paragraph{Decomposition strategies compared directly.} Several studies compare a single call with a decomposed
alternative, with heterogeneous tasks and results. In LLM-IE, a sentence-by-sentence extractor achieved higher
F1 than a single whole-document call on three English clinical corpora (for example 0.78 versus 0.44 for event
extraction), mostly through recall, at about twice the GPU time per note; the authors recommend it for ``dense''
tasks and the single call for ``sparse'' ones~\cite{llmie2025}. Fornasiere et al.\ compared zero-shot, few-shot
and sequential prompting with a single 7B model on discharge summaries; sequential prompting did not outperform
few-shot prompting for timeline extraction, and extraction of full medication details remained much harder
than extraction of names~\cite{fornasiere2024medical}. Notably, they supplied the note with numbered rows and
asked the model for the row supporting each extraction, the same principle as the line identifiers used here,
but did not evaluate the correctness of those references. In a preprint benchmark of 15 open-weight models on
Dutch, English and Czech pathology and radiology reports, a manually defined prompt graph of smaller sequential
subtasks gave the largest mean gain over zero-shot prompting ($+12.7\%$), at 37.6~s per report versus 4.5~s for
few-shot prompting; all other strategies requested the whole schema in one JSON call, with at most 22 fields per
template~\cite{spaanderman2025structured}. Finally, CLEAR decomposed the input rather than the output: for each
target variable it retrieved entity-anchored windows and called the model once per window, obtaining higher F1
than full-note input with 71\% fewer input tokens and inference time reduced by 72\%~\cite{clinicalentityretrieval2024}.

\paragraph{Reporting.} A scoping review of 34 studies on information extraction from radiology reports, which
included no generative or prompt-based approaches, found that performance could not be compared across studies
because of differences in data, number of concepts and metrics; 68\% of studies did not mention annotation
guidelines, inter-annotator agreement values were given by 16 of 34, and code was available for 29\%
\cite{scoping2024radiology}. TRIPOD-LLM now asks studies that evaluate LLMs to report prompt development and
its data, inference settings, annotation guidelines, the proportion of multiply annotated data,
inter-annotator agreement, annotator background and compute~\cite{gallifant2025tripodllm}.

\paragraph{Gap.} Across these studies, decomposition has been applied to the input (sentences, retrieved
windows), to the reasoning (segmentation then standardisation, prompt graphs) or to a review pass, and its
reported benefits come with higher cost (\cref{tab:sota}). None of the studies reviewed here compares, on the
same patients and models, a single call with an output-side decomposition by schema group for a form of the
size used here: the largest schema requested in a single call among them has 22 fields, against 199 in this
study. Evidence grounding was either not evaluated or limited to retrieval, and line-numbered evidence, where
used, was not verified. The present study targets that gap with a paired design, verifiable line-level
evidence, full execution logs and open-weight models on local infrastructure.

\subsection{Objectives and hypotheses}\label{sec:objectives}

\paragraph{Research question.} In a fixed cohort of Spanish discharge summaries of patients admitted to an
adult ICU, how do three ways of requesting the same 199-variable form compare when the note, the rules, the
output contract, the validator and the model are held constant: (A) a single call for the whole form;
(B) 18 independent calls, one per group of related variables; and (C) the same 18 calls sharing an
evidence-grounded model of the case built by a preliminary call? The comparison covers structural validity,
evidence fidelity, clinical coherence across blocks, agreement with human annotators and cost, and asks
whether any benefit of the shared case model depends on how much inference a variable requires and on the
quality of the note.

\paragraph{Pre-specified hypotheses.} The following are hypotheses, not findings. H1--H4 were specified
before the pilot for Arms A and B. After the pilot, and before any output was compared with the human
reference, H1 was refined to separate field-level failures from relation violations between groups, H1--H4
were extended to Arm~C, and H5 and H6 were added (Section~\ref{sec:separation}); the original wording is
kept in the protocol history.
\begin{description}
\item[H1 (structure).] Arms B and C have fewer truncated or unparseable outputs and fewer field-level
failures than Arm A. Arm B has more cross-group relation violations than Arm A, because its groups are
answered independently, and Arm C has fewer than Arm B.
\item[H2 (evidence).] Arms B and C have a higher proportion of positive and uncertain answers supported by
valid, relevant evidence lines than Arm A. Because all arms cite by line identifier, literal correspondence
holds by construction; H2 concerns identifier validity, coverage and clinician-judged support.
\item[H3 (agreement, primary).] Arm B is non-inferior to Arm A in agreement with the pre-specified human
reference on the 174 boolean variables of eligible summaries; Arm C is non-inferior to Arm A.
\item[H4 (cost).] Decomposition increases calls, input tokens and total computation per case (A $<$ B $<$ C
in calls and input tokens); latency depends on the inference engine and is reported, not tested.
\item[H5 (clinical coherence).] Arm C violates fewer pre-specified cross-block coherence rules per case than
Arm B.
\item[H6 (inference and note quality).] Arm C improves agreement over Arm B in the blocks that require
inference across the episode (infections, organ failure, complications and ICU discharge), is non-inferior
in medical history, and the improvement is larger in low-quality notes.
\end{description}
A single-case pilot is used only to test the pipeline, the instrumentation and the error handling;
it does not contribute to the estimation of any effect.

\begin{landscape}
\begin{table}[p]
\centering\scriptsize
\setlength{\tabcolsep}{3pt}
\caption{\textbf{Studies used to position this work.} Information as reported by each paper; ``not reported''
means the paper does not state it. IAA, inter-annotator agreement; Source:
\texttt{referencias/INDEX.md} (full texts in \texttt{referencias/md}).}
\label{tab:sota}
\begin{tabularx}{\linewidth}{>{\raggedright\arraybackslash}p{1.7cm}>{\raggedright\arraybackslash}p{2.1cm}>{\raggedright\arraybackslash}p{2.3cm}YY>{\raggedright\arraybackslash}p{2.3cm}>{\raggedright\arraybackslash}p{2.3cm}Y>{\raggedright\arraybackslash}p{1.3cm}Y}
\toprule
Study & Language; notes & Annotators (IAA) & Output schema & Prompt strategy & Models & Evidence validation & Main metric (example) & Code & Main limitation \\
\midrule
Agrawal 2022~\cite{agrawal2022fewshot} & English; clinical snippets, abstracts & 2 authors, adjudicated (no $\kappa$ for new sets) & Free text or lists mapped by resolvers; up to 5 medication attributes & One-shot guided prompts; one item, or medications with status, per prompt & GPT-3 (API) & Resolvers drop tokens absent from input & Medication extraction recall 0.90, precision 0.92 & Data released; code not reported & Hard to match an exact schema; non-empty answer bias \\
Wiest 2024~\cite{privacy2024structured} & English; MIMIC-IV history sections (500) & 3 blinded experts, consensus (no coefficient) & 5 binary features, JSON with excerpt & Fixed prompt with grammar, all features in one JSON (calls per report not stated); zero/one-shot, definitions, excerpt-first & Llama-2 7B--70B, local & Excerpt requested, not checked; hallucinated excerpt reported & Accuracy 0.93--0.96 (70B, zero-shot with definitions) & Yes & English only; lower recall for implicit features \\
Hein 2025~\cite{iterativerefine2025} & Not stated (English excerpts); kidney-tumour pathology reports & Data scientist, statistician, pathologist for ambiguity (no IAA) & JSON to long table; controlled vocabularies & Segmentation then standardisation, per feature or entity class; reasoning passed forward & GPT-4o (API); Llama 3.3 70B, Qwen2.5 72B local & None automated; human discrepancy review & Major error 0.99\%; internal macro-F1 0.99 & Yes & No IAA; no held-out benchmark over all fields \\
Fornasiere 2024~\cite{fornasiere2024medical} & English (+ synthetic Italian); discharge summaries & Challenge gold standards (no new IAA) & Medications (5 fields); dated events & Zero-/few-shot single call; 3-turn sequential; whole vs chunked & Mistral 7B & Row numbers requested, not evaluated & Name F1 up to 0.683; timeline F1 0.790 few-shot vs 0.660 sequential & Pledged, no URL & One model; small datasets \\
Hsu 2025 (LLM-IE)~\cite{llmie2025} & Not stated; i2b2/n2c2 clinical notes & Challenge gold standards & Frames with spans, attributes, relations & Whole document vs per sentence vs review pass; 8-shot & Llama-3.1-70B, self-hosted & Spans in frames; no dedicated verification & Event F1 0.78 per sentence vs 0.44 whole document; $\sim$2$\times$ GPU time & Yes & One model benchmarked \\
Lopez 2025 (CLEAR)~\cite{clinicalentityretrieval2024} & Not stated; clinical notes, radiology reports & 5 physicians, 1 student ($\kappa$ 0.86 and 0.93) & 1 ternary label per variable; 18 variables & Per-variable entity retrieval; one call per retrieved chunk; 5-shot & 5 local models (up to 70B); GPT-4 (Azure) & Retrieval-level only & F1 0.90 vs 0.79 full note; 71\% fewer input tokens & Yes & Multi-chunk prompts and quantisation not tested \\
Spaanderman 2025 (preprint)~\cite{spaanderman2025structured} & Dutch, English, Czech; pathology and radiology reports & Varies by use case (agreement 0.68--0.88 where measured) & Flat JSON, 4--22 fields per template & Zero/one/few-shot, CoT, self-consistency (single call) vs prompt graph & 15 open-weight, 1.5B--685B, local & None evaluated & Prompt graph $+12.7\%$ vs zero-shot; 37.6 vs 4.5~s per report & Yes & No quantisation study; mostly pathology \\
Chen 2026 (two-stage)~\cite{psych2stage2026} & Mostly English; psychiatric discharge summaries (500) & Psychiatrist and NLP researcher, adjudicated ($\kappa$ 0.78) & Sentence events and times; 4 chart-level features & Per-sentence extraction, then one chart-level call on the intermediate list; vs direct and joint & 4 open 8B models, LoRA, local & Source sentence kept; not verified & Onset accuracy 0.73--0.77 vs 0.56--0.65 direct; no gain for counts & Not reported & Single hospital; 4 features \\
Aal Abdulsalam 2026 (multi-agent)~\cite{cancerregistrymas2026} & Not stated; pathology and medical reports (628 evaluated) & Pathologists, senior adjudication ($\kappa$ 0.92) & 6 registry fields, closed labels, evidence quote & Chunking, retrieval, one agent per field, review, aggregation; vs direct prompting & Llama 3.3 70B, local & Quote requested, not checked & Grade F1 0.71$\to$0.78; T and M worse; mean F1 0.71 vs 0.74 & Yes & No ablation \\
Pathak 2026 (preprint)~\cite{promptplanextract2026} & English; lung pathology reports (test 53) & 3 annotators and expert ($\kappa$ 0.81) & 13 CAP fields, JSON & Section mapper, planner metadata injected into parallel field executors, compiler & 5 open models & None & Micro-F1 0.83--0.89 vs 0.96 supervised & Not reported & No single-call baseline or ablation \\
\midrule
This study & Spanish; ICU discharge summaries (50) & 6 assigned, 4 submitting; $\kappa$ 0.76 on current cut & 199 typed variables with hierarchy and exclusions & 1 call vs 18 schema-group calls vs 18 calls sharing a validated case model & Gemma 4 31B, Llama 3.3 70B, Qwen 3.6 35B-A3B, local & Line IDs validated; text copied by software & \pending{paired differences} & Versioned & Single centre; reference incomplete \\
\bottomrule
\end{tabularx}
\end{table}
\end{landscape}
